\documentclass[10pt]{article}
\usepackage[T1]{fontenc}
\usepackage[utf8]{inputenc}
\usepackage[french]{babel}
\usepackage{amsmath,amssymb,mathrsfs,mathpazo}
\usepackage{geometry}
\geometry{a4paper,margin=2cm}
\usepackage{graphicx,array,tabularx,tabularray,colortbl,multirow,multicol,booktabs,enumitem,pifont,ulem,fancybox,diagbox,cancel}
\usepackage{pstricks,pstricks-add,pst-plot,pst-tree,pst-node,pst-text,pst-func,pst-math,pst-bspline,pst-3dplot,pst-eucl,pst-solides3d}
\let\pscalculate\undefined
\usepackage{pst-bezier}
\usepackage{tikz,pgfplots,tkz-tab}
\usetikzlibrary{arrows,arrows.meta,patterns,calc,positioning,trees,decorations.markings,shapes.geometric,shapes.misc,shapes.arrows,intersections,decorations.pathmorphing,decorations.pathreplacing,backgrounds,shadows,babel}
\pgfplotsset{compat=1.18}
\usepgfplotslibrary{fillbetween}
\usepackage[np]{numprint}
\usepackage{listings,makecell,pas-tableur,twoopt}
\let\barre\undefined
\usepackage[pstricks]{tablvar}
\definecolor{darkgreen}{RGB}{0,100,0}
\providecommand{\red}{\color{red}}
\providecommand{\blue}{\color{blue}}
\providecommand{\green}{\color{green}}
\providecommand{\black}{\color{black}}
\providecommand{\white}{\color{white}}
\providecommand{\gray}{\color{gray}}
\providecommand{\cyan}{\color{cyan}}
\providecommand{\magenta}{\color{magenta}}
\providecommand{\yellow}{\color{yellow}}
\providecommand{\violet}{\color{violet}}
\providecommand{\orange}{\color{orange}}
\providecommand{\euro}{\texteuro}
\providecommand{\EUR}{\texteuro}
\providecommand{\ieme}{\textsuperscript{e}}
\providecommand{\degres}{\ensuremath{^\circ}}
\providecommand{\degree}{\ensuremath{^\circ}}
\providecommand{\vv}[1]{\overrightarrow{#1}}
\providecommand{\Coord}[1]{#1}
\providecommand{\decosix}{}
\providecommand{\starredbullet}{\ensuremath{\bullet}}
\pagestyle{empty}
\setlength{\parindent}{0pt}
\usepackage{adjustbox}\DeclareUnicodeCharacter{2236}{:}\DeclareUnicodeCharacter{2190}{\ensuremath{\leftarrow}}
\providecommand{\R}{\mathbb{R}}
\providecommand{\N}{\mathbb{N}}
\providecommand{\Z}{\mathbb{Z}}
\providecommand{\Q}{\mathbb{Q}}
\providecommand{\D}{\mathbb{D}}
\providecommand{\C}{\mathbb{C}}
\providecommand{\e}{{\mathrm{e}}}
\providecommand{\vect}[1]{\overrightarrow{#1}}
\providecommand{\vectt}[1]{\overrightarrow{\mathrm{#1}}}
\providecommand{\vvt}[1]{\overrightarrow{\mathrm{#1}}}
\providecommand{\np}[2][]{{#2~#1}}
\providecommand{\vv}[1]{{\overrightarrow{#1}}}
\renewcommand{\d}{{\mathrm{d}}}
\providecommand{\degres}{{^\circ}}
\providecommand{\degre}{{^\circ}}
\providecommand{\blue}{\color{blue}}
\providecommand{\red}{\color{red}}
\providecommand{\enskip}{\quad}
\providecommand{\llbracket}{\lbrack\!\lbrack}
\providecommand{\rrbracket}{\rbrack\!\rbrack}
\providecommand{\euro}{{\text{€}}}
\providecommand{\no}{{\text{n°}}}
\providecommand{\interval}[2]{{[#1;#2]}}
\providecommand{\intervFO}[2]{{[#1;#2[}}
\providecommand{\dotfill}{\ldots}

\providecommand{\ds}{\displaystyle}
\providecommand{\textding}[1]{#1}
\providecommand{\Oij}{\(\mathrm{O};\vec{\imath},\vec{\jmath}\)}
\providecommand{\Oijk}{\(\mathrm{O};\vec{\imath},\vec{\jmath},\vec{k}\)}
\providecommand{\Ouv}{\(\mathrm{O};\vec{u},\vec{v}\)}
\providecommand{\intervFF}[2]{\ensuremath{[#1;#2]}}
\renewcommand{\euro}{\ifmmode\text{\texteuro}\else\texteuro\fi}
\DeclareUnicodeCharacter{00B0}{\ifmmode^{\circ}\else\textdegree\fi}
\DeclareUnicodeCharacter{20AC}{\ifmmode\text{\texteuro}\else\texteuro\fi}
% Énoncés : Éducation nationale. Transcription et corrigés : APMEP, Denis Vergès et contributeurs. Découpage et classement : CoopMaths, Sébastien Lozano, Éric Elter et contributeurs. Adaptation de la présentation : Maths Explorer.
% CC BY-SA : https://coopmaths.fr/www/about/
% https://coopmaths.fr/alea/static/e3c/2020/tex/e3c_2020_00_sujet43_3_cor.tex
\begin{document}
\noindent\textit{APMEP / CoopMaths. CC BY-SA. Présentation adaptée par Maths Explorer.}\par\medskip
\subsection*{1.}

On sait que :  et , donc :
\[
\begin{cases}
A(-1\,;\,3) \in (AB) \\
B(5\,;\,0) \in (AB)
\end{cases}
\iff
\begin{cases}
3 = -a + b \\
0 = 5a + b
\end{cases}
\Rightarrow -3 = 6a \iff a = -\dfrac{1}{2},
\]
puis \( b = a + 3 = -\dfrac{1}{2} + 3 = \dfrac{5}{2} \).

On a donc :
\[
M(x\,;\,y) \in (AB) \iff y = -\dfrac{1}{2}x + \dfrac{5}{2}.
\]

\subsection*{2.}

\begin{align*}
&M(x\,;\,y) \in \mathcal{D} \\
\iff &\overrightarrow{CM} \cdot \vec{n} = 0 \\
\iff &- (x - 9) + 3(y - 3) = 0 \\
\iff& -x + 3y = 0.
\end{align*}

\subsection*{3.}

\( \mathcal{D} \) a pour vecteur directeur \( \vec{d} \begin{pmatrix} -3 \\ -1 \end{pmatrix} \) et \( (AB) \) a pour vecteur directeur \( \overrightarrow{AB} \begin{pmatrix} 6 \\ -3 \end{pmatrix} \).

Or, \( \det\left(\vec{d}, \overrightarrow{AB}\right) = 9 + 6 = 15 \neq 0 \) : les droites \( \mathcal{D} \) et \( (AB) \) ne sont pas parallèles.

\subsection*{4.}

\[
\vec{d} \cdot \overrightarrow{AB} = -3 \times 6 + (-1) \times (-3) = -18 + 3 = -15 \neq 0,
\]
ces vecteurs ne sont pas orthogonaux, donc les droites \( \mathcal{D} \) et \( (AB) \) ne sont pas perpendiculaires.

\subsection*{5.}

On a :
\[AC^2 = (9 - (-1))^2 + (3 - 3)^2 = 10^2,
\]
donc \( AC = 10 \).

D'après le théorème d'Al-Kashi :
\begin{align*}
AC^2 &= AE^2 + EC^2 - 2 \times AE \times EC \times \cos \widehat{AEC} \\
100 &= \left(2\sqrt{5}\right)^2 + \left(2\sqrt{10}\right)^2 - 2 \times 2\sqrt{5} \times 2\sqrt{10} \times \cos \widehat{AEC} \\
100 &= 20 + 40 - 8\sqrt{50} \cos \widehat{AEC} \\
8\sqrt{50} \cos \widehat{AEC} &= -40,
\end{align*}

d'où l'on exprime :
\begin{align*}
\cos \widehat{AEC} &= -\dfrac{5}{\sqrt{50}} \\
&= -\dfrac{5}{\sqrt{5^2 \times 2}} \\
&= -\dfrac{5}{5\sqrt{2}} \\
&= -\dfrac{5}{5\sqrt{2}} \\
&= -\dfrac{1}{\sqrt{2}} \\
&= -\dfrac{\sqrt{2}}{2}.
\end{align*}

On en déduit que \( \widehat{AEC} = \dfrac{3\pi}{4} \).

\end{document}
